\documentclass[10pt,a4paper]{amsart}
\usepackage[margin=19mm]{geometry}
\usepackage{amsmath,amssymb,mathtools}
\usepackage[scr=rsfs,cal=euler]{mathalfa}
\usepackage{xfrac}
\usepackage[unicode,hidelinks,bookmarks=false]{hyperref}
\newcommand{\ie}{i.e.\ }
\newcommand{\cf}{cf.\ }
\newcommand{\resp}{respectively\ }
\newcommand{\Subsection}{\subsection}
\providecommand{\nobreakdash}{\nobreak\hbox{-}}
\setlength{\emergencystretch}{2em}
\multlinegap0pt
\usepackage{amssymb}
\usepackage[shortlabels]{enumitem}
\usepackage{xcolor}
\usepackage{mathtools}
\usepackage{tensor}
\usepackage{tikz}
\usepackage{cancel}
\newtheorem{lemma}{Lemma}
\newcommand{\ud}{\,\mathrm{d}}
\title{Quantitative stability for the Trudinger--Moser inequality}
\author{Jo\~ao Henrique Andrade, Jos\'e Francisco de Oliveira, Jo\~ao Marcos do \`O, Abiel Costa Macedo, Jesse Ratzkin}
\date{Source: arXiv:2605.03668v1 (CC BY 4.0)}
\begin{document}
\maketitle

\noindent\emph{Source and licence.} Quantitative stability for the Trudinger--Moser inequality. Version arXiv:2605.03668v1 (CC BY 4.0), distributed by arXiv under the Creative Commons Attribution 4.0 International licence, http://creativecommons.org/licenses/by/4.0/, as recorded in the arXiv licence metadata for this exact version and shown on the abstract page https://arxiv.org/abs/2605.03668v1. Full licence notice: \texttt{licenses/arxiv-2605.03668-CC-BY-4.0.txt}.\par\smallskip

\noindent\textit{Editorial scope.} Counted unit: the article's lemma lm:taylor\_expansion (source0.tex lines 1113--1127). The article's complete original proof of it is delivered with it (source0.tex lines 1129--1334, one \texttt{proof} environment of 206 lines). No mathematics is written, repaired or re-derived: the statement and the proof are the article's own lines, in the article's own notation, and no retained line is substituted or re-flowed.  Retained uncounted as the source-local context the proof needs: lines 522--527; lines 960--962.  Nothing was omitted from any retained span, so the package carries no omission record.  No retained line contains a citation, so the wrapper carries no bibliography and no bibliography entry.  No retained line contains a figure environment, an image-inclusion command or drawing code, so no image is delivered and the package declares no asset dependency.  Editorial formatting only, in addition: an AMS \texttt{amsart} preamble that restates the article's own declarations and macros used by the retained lines, namely the environments \texttt{lemma} on separate counters, together with the standard packages those lines need.  The retained environments print in this document as Lemma 1 in order of appearance, whereas the article prints the same items with the numbering of its own full document.  The complete native source of the article as distributed in the arXiv e-print for this exact version is bundled unchanged as \texttt{sources/199806/source0.tex}.
\begin{equation}\label{eq:TM_BVP}\tag{$\mathcal{TM}_{\alpha,\Omega}$}
		 \begin{cases}
			-\Delta v = \Lambda  e^{\alpha u_{*}^2} v & \mbox{in} \quad \Omega\\
			v=0 & \mbox{on} \quad\partial \Omega.
		\end{cases} 
	\end{equation}

\begin{equation}\label{eq:first_TM_eigenvalue}
 \Lambda_1({\alpha})=\frac{1}{\int_{\Omega} u_*^2 e^{\alpha u_*^2} \ud x}
\end{equation}

\begin{lemma}[Expansion near optimizers]\label{lm:taylor_expansion}
Let $\Omega\subset\mathbf{R}^2$ be a smooth, bounded domain.
For any $u_*\in \mathcal{M}_\alpha(\Omega)\cap \mathcal{S}_1$ and $u\in \mathcal{S}_1$ one has 
\begin{equation} \label{eq:taylor_expansion} 
\mathscr{E}_{\alpha,\Omega} (u_*) - \mathscr{E}_{\alpha,\Omega} (u) 
= \frac{\alpha}{\Lambda_1(\alpha)} 
		\| \nabla w\|_{L^2(\Omega)}^2 - \alpha\mathcal{I}^*_{\alpha,\Omega}(w) 
        + \mathcal{O}(\|\nabla w\|_{L^2(\Omega)}^3) \quad {\rm as} \quad \|\nabla w\|_{L^{2}(\Omega)}\to0,
\end{equation} 
where $w=u-u_*$ and
\begin{equation}\label{eq:integral_error}
  \mathcal{I}^*_{\alpha,\Omega}(w)= \int_\Omega w^2(1+2\alpha u_*^2)
		e^{\alpha u_*^2} \ud x.  
\end{equation}
\end{lemma} 

\begin{proof}
Initially, a direct computation implies that the following expansion holds: 
\begin{align} \label{expansion1} 
\nonumber
\mathscr{E}_{\alpha,\Omega} (u_* + w)
&= \int_\Omega e^{\alpha u_*^2} \sum_{\ell =0}^\infty \frac{\alpha^\ell}{\ell !}
(2u_* w + w^2)^\ell \ud x \\ \nonumber 
& = \int_\Omega e^{\alpha u_*^2} \ud x + \alpha \int_\Omega e^{\alpha u_*^2} 
(2u_* w + w^2) \ud x + \frac{\alpha^2}{2} \int_\Omega e^{\alpha u_*^2} (4u_*^2 w^2 
+ 4u_* w^3 + w^4) \ud x\\\nonumber
&+ \sum_{\ell=3}^\infty \frac{\alpha^\ell}{\ell !} \int_\Omega e^{\alpha u_*^2} 
(2u_* w + w^2)^\ell \ud x\\
&:=\mathscr{E}_{\alpha,\Omega} (u_*)+\alpha\mathcal{I}^*_{\alpha,\Omega}(w)+T_1(w)+T_2(w)+T_3(w),
\end{align} 
where $\mathcal{I}^*_{\alpha,\Omega}$ is defined in \eqref{eq:integral_error} and
\[
T_1(w) = 2\alpha \int_\Omega u_* w e^{\alpha u_*^2} \ud x,
\]
\[
T_2(w) = \int_{\Omega} \left (2\alpha^2 u_* w^3 + 
\frac{\alpha^2}{2} w^4 \right ) e^{\alpha u_*^2} \ud x, 
\]
and
\[
T_3(w) = \sum_{\ell=3}^\infty \frac{\alpha^\ell}{\ell !} \int_\Omega e^{\alpha u_*^2} 
(2u_* w + w^2)^\ell \ud x. 
\]
Then
\[\mathscr{E}_{\alpha,\Omega} (u_*) - \mathscr{E}_{\alpha,\Omega} (u) =-T_1(w)-\alpha\mathcal{I}^*_{\alpha,\Omega}(w)-T_2(w)-T_3(w).\]
We analyze each of these terms individually.
We start with the first term, for which we use the first eigenvalue of the boundary problem \eqref{eq:TM_BVP}. 

\noindent{\bf Claim 1:} $T_1(w)=-\frac{\alpha}{\Lambda_1(\alpha)}\|\nabla w\|_{L^2(\Omega)}^2$.

\noindent First, observe
that since $u_*,u\in \mathcal{S}_1$, one has
\[1 = \int_\Omega |\nabla u|^2 \ud x = \int_\Omega |\nabla u_* + \nabla w|^2 \ud x =
\int_\Omega |\nabla u_*|^2 \ud x + 2 \int_\Omega \langle \nabla u_*, \nabla w \rangle
\ud x + \int_\Omega |\nabla w|^2 \ud x,\]
which, by rearranging, implies  
\begin{equation} \label{first_order1}  
\int_\Omega |\nabla w|^2 \ud x = -2 \int_\Omega \langle \nabla u_*, \nabla w
\rangle \ud x.
\end{equation} 
Second, since $u_*\in\mathcal{M}_\alpha(\Omega)$ is an extremal of the functional $\mathscr{E}_{\alpha,\Omega}$, it solves the 
boundary value problem \eqref{eq:TM_BVP} with Lagrange multiplier $\Lambda_1(\alpha)>0$ defined in \eqref{eq:first_TM_eigenvalue}, namely 
\begin{equation}\label{euler_lag_eqn}
		 \begin{cases}
			-\Delta u_*  = \Lambda_1(\alpha)  e^{\alpha u_{*}^2} u_* & \mbox{in} \quad \Omega\\
			u_*=0 & \mbox{on} \quad\partial \Omega.
		\end{cases} 
	\end{equation}
As a consequence of \eqref{first_order1} and 
\eqref{euler_lag_eqn}, we get
\begin{eqnarray} \label{first_order2} 
T_1(w) =  \frac{2\alpha}{\Lambda_1(\alpha)} \int_\Omega \langle \nabla u_*, \nabla w \rangle \ud x = -\frac{\alpha}{\Lambda_1(\alpha)} \int_\Omega |\nabla w|^2\ud x = -\frac{\alpha}{\Lambda_1(\alpha)} 
\| \nabla w \|_{L^2(\Omega)}^2. 
\end{eqnarray} 
This proves the first estimate.

Before analyzing the terms $T_2(w)$ and $T_3(w)$,
let us perform some preliminary computations. 
First, observe that the convexity of the function $t \mapsto t^\ell$ implies that 
for each $a,b >0$ we have 
\begin{equation}\label{convex_power}
\left ( \frac{a + b}{2} \right )^\ell \leq \frac{a^\ell + b^\ell}{2} \quad {\rm or} \quad
(a+b)^\ell \leq 2^{\ell-1} (a^\ell + b^\ell ).
\end{equation} 
Next, we use the Trudinger--Moser and H\"older 
inequalities to bound integrals of powers of $w$. 
More precisely, we have 
    \begin{align*} 
    {\rm TM}(\alpha,\Omega) \geq \int_\Omega \exp \left (\frac{\alpha w^2}{\| \nabla w\|_{L^2(\Omega)}^2} \right ) \ud x 
    & = \int_\Omega \sum_{j = 0}^\infty \frac{\alpha^j}{j ! 
    \| \nabla w \|_{L^2(\Omega)}^{2j} }w^{2j}\ud x\\
    &= \sum_{j = 0}^\infty \frac{\alpha^j}
    {j ! \| \nabla w\|_{L^2(\Omega)}^{2j}} 
    \int_\Omega w^{2j} \ud x \\ 
    & \geq \frac{\alpha^\ell}{\ell! \| \nabla w \|_{L^2(\Omega)} ^{2\ell}} 
    \int_\Omega w^{2\ell} \ud x\\
    &= \frac{\alpha^{\ell} \| w\|_{L^{2\ell}(\Omega)}^{2\ell}}
    {\ell! \| \nabla w\|_{L^2(\Omega)}^{2\ell} }
    \end{align*} 
    for any specific $\ell \in \mathbf{N}$. This is justified by the fact that $w^2 \geq 0$, 
    so each term in the sum is non-negative. We can rearrange this inequality to read 
    \begin{equation} \label{higher_order1}
    \int_\Omega |w|^{2\ell}\ud x \leq {\rm TM}(\alpha,\Omega)\frac{\ell!}{\alpha^\ell}
    \| \nabla w \|_{L^2(\Omega)}^{2\ell} . \end{equation}
    We can similarly control the integral of $w^{2\ell-1}$ using H\"older's inequality as follows: 
  \begin{align} \label{higher_order2} 
    \| w \|_{L^{2\ell-1}(\Omega)}^{2\ell -1} & =  \| w^{2\ell -1} \|_{L^1(\Omega)}    \leq \| \chi_{\Omega} \|_{L^{2\ell}(\Omega)} 
    \| w^{2\ell -1} \|_{L^{\frac{2\ell}{2\ell-1}}(\Omega)} \\ \nonumber 
    & = |\Omega|^{\frac{1}{2\ell}} \left ( \int_\Omega w^{2\ell} \ud x\right )^{\frac{2\ell-1}{2\ell}} 
    = |\Omega|^{\frac{1}{2\ell}} \left ( 
    \| w \|_{L^{2\ell}(\Omega)}^{2\ell} \right )^{\frac{2\ell-1}{2\ell}} \\ \nonumber 
    & \leq {\rm TM}(\alpha,\Omega)^{\frac{2\ell-1}{2\ell}} \frac{(\ell !)^{\frac{2\ell-1}{2\ell}}}
    {\alpha^{\frac{2\ell-1}{2}}} \| \nabla w\|_{L^2(\Omega)}^{2\ell -1}.
    \end{align} 

    We can proceed to estimate the second term. 
    
    \noindent{\bf Claim 2:} $T_2(w)=\mathcal{O}(\|\nabla w\|^3_{L^2(\Omega)})$ as $\|\nabla w\|_{L^2(\Omega)}\to0$.

    \noindent
    As a matter of fact, using \eqref{higher_order1} we see 
    \begin{equation} \label{higher_order8} 
    \left | \frac{\alpha^2}{2} \int_\Omega e^{\alpha u_*^2} w^4 \ud x \right | \leq 
    {\rm TM}(\alpha,\Omega) |\Omega |e^{\alpha M^2} \| \nabla w \|_{L^2(\Omega)}^4 \lesssim_{\alpha, \Omega, M}\| \nabla w\|_{L^2(\Omega)}^4.\end{equation} 
    Additionally, by means of \eqref{higher_order2}, we get
    \begin{equation} \label{higher_order9} 
    \left | 2\alpha^2 \int_\Omega u_*w^3 e^{\alpha u_*^2} \ud x \right | \leq 
    2^{7/4} \alpha ^{1/4} {\rm TM}(\alpha,\Omega)^{3/4} M e^{\alpha M^2} 
    \| \nabla w\|_{L^2(\Omega)}^3\lesssim_{\alpha, \Omega, M}\| \nabla w\|_{L^2(\Omega)}^3.\end{equation}
    This concludes the proof of the claim. 


    At last, we can estimate the third term, which is slightly more technical.
    
    \noindent{\bf Claim 3:} $T_3(w)=\mathcal{O}(\|\nabla w\|^3_{L^2(\Omega)})$ as $\|\nabla w\|_{L^2(\Omega)}\to0$.

    \noindent Now, let $M = \max \{ 1, \| u_*\|_{L^\infty (\Omega)} \}$. Using \eqref{convex_power}
    we can bound $T_3(w)$ by the following three series: 
    \begin{align} \label{higher_order3} 
    |T_3(w)| & =  \left | \sum_{\ell=3}^\infty \frac{\alpha^\ell}{\ell !} \int_\Omega 
    e^{\alpha u_*^2} (2u_* w + w^2)^{\ell} \ud x \right | \\ \nonumber 
    & \leq  e^{\alpha M^2} \sum_{\ell=3}^\infty  \frac{\alpha^\ell 2^{\ell-1}}{\ell !} 
    \int_\Omega (2^\ell u_*^\ell |w|^\ell + |w|^{2\ell} ) \ud x \\ \nonumber 
    & \leq e^{\alpha M^2} \left [ \sum_{\ell=2}^\infty \frac{\alpha^{2\ell-1} 2^{4\ell-3} M^{2\ell-1} }
    {(2\ell-1) !} \int_\Omega |w|^{2\ell-1} \ud x + \sum_{\ell=2}^\infty \frac{\alpha^{2\ell} 2^{4\ell-1}
    M^{2\ell} }{(2\ell)!} \int_\Omega |w|^{2\ell} \ud x  \right ] \\ \nonumber 
    &+ e^{\alpha M^2} \sum_{\ell=3}^\infty \frac{\alpha^\ell 2^{\ell-1}}{\ell !}
    \int_\Omega |w|^{2\ell} \ud x\\ \nonumber 
    &:= e^{\alpha M^2} [S_1(w)+ S_2(w)+S_3(w)]. 
    \end{align} 
    Here we set
    \[
    S_1(w):= \sum_{\ell=2}^\infty \frac{\alpha^{2\ell-1} 2^{4\ell-3} M^{2\ell-1} }
    {(2\ell-1) !} \int_\Omega |w|^{2\ell-1} \ud x,
    \]
    \[
    S_2(w):=\sum_{\ell = 2}^\infty \frac{\alpha^{2\ell} 2^{4\ell-1}M^{2\ell} }{(2\ell)!}
    \int_\Omega |w|^{2\ell} \ud x,
    \]
    and
    \[
    S_3(w):=\sum_{\ell=3}^\infty \frac{\alpha^\ell 2^{\ell-1}}{\ell !}
    \int_\Omega |w|^{2\ell} \ud x.
    \]
    \noindent Next, we estimate each sum individually. First, we use \eqref{higher_order1} to bound the 
    last sum by a geometric series, so long as 
    $\| \nabla w \|_{L^2(\Omega)}^2 < \frac{1}{2} .$
    More precisely, one has
    \begin{align*} 
    S_3(w) & = \sum_{\ell=3}^\infty \frac{\alpha^\ell 2^{\ell-1}}{\ell !}
    \int_\Omega |w|^{2\ell} \ud x \\ \nonumber 
    & \leq {\rm TM}(\alpha,\Omega) \frac{|\Omega |}{2} \sum_{\ell=3}^\infty(2\| 
    \nabla w\|_{L^2(\Omega)}^2)^\ell \\ \nonumber 
    & = {\rm TM}(\alpha,\Omega)\frac{|\Omega|}{2} \left ( \frac{1}{1-2\|\nabla w\|_{L^2(\Omega)}^2}
    - 1 - 2\| \nabla w\|_{L^2(\Omega)}^2 - 4\| \nabla w\|_{L^2(\Omega)}^4 \right ) \\ \nonumber 
    & = {\rm TM}(\alpha,\Omega)\left(\frac{4\| \nabla w \|_{L^2(\Omega)}^6}{1-2\| \nabla w \|_{L^2(\Omega)}^2}\right). 
    \end{align*} 
    From this, we conclude
    \begin{equation} \label{higher_order4}
    |S_3(w)| \lesssim_{\alpha,\Omega} \| \nabla w \|_{L^2(\Omega)}^4.
    \end{equation}
    We can also use \eqref{higher_order1} to bound the second sum when $0<\alpha\ll1$. Also, when $\| \nabla w \|_{L^2(\Omega)} < 1$, we have  
    \begin{eqnarray*} 
    S_2(w) & = & \sum_{\ell = 2}^\infty \frac{\alpha^{2\ell} 2^{4\ell-1}M^{2\ell} }{(2\ell)!} 
    \int_\Omega |w|^{2\ell} \ud x \\  
    & \leq & {\rm TM}(\alpha,\Omega) |\Omega| \sum_{\ell=3}^\infty \frac{\alpha^\ell 2^{4\ell-1} \ell! M^{2\ell}}
    {(2\ell)!} \| \nabla w\|_{L^2(\Omega)}^{2\ell} \\  
    & \leq & {\rm TM}(\alpha,\Omega) \| \nabla w\|_{L^2(\Omega)}^4 \sum_{\ell=2}^\infty 
    \frac{\alpha^\ell 2^{4\ell-1} \ell! M^{2\ell}}{(2\ell)!}. 
    \end{eqnarray*} 
    One can show, using the ratio test, that this last sum converges; and so
    \begin{equation} \label{higher_order5} 
    |S_2(w)| \lesssim_{\alpha,\Omega,M} \| \nabla w \|_{L^2(\Omega)}^4.
    \end{equation}
    At last, we use \eqref{higher_order2} to bound $S_1(w)$. 
    In this case, we have 
    \begin{align*} 
    S_1(w) & = \sum_{\ell=2}^\infty \frac{\alpha^{2\ell-1} 2^{4\ell-3} M^{2\ell-1} }
    {(2\ell-1) !} \int_\Omega |w|^{2\ell-1} \ud x \\ 
    & \leq |\Omega| \sum_{\ell=2}^\infty \frac{\alpha^{\frac{2\ell-2}{2}} 2^{4\ell-3} 
    M^{2\ell-1} {\rm TM}(\alpha,\Omega)^{\frac{2\ell-1}{2\ell}} (\ell!)^{\frac{2\ell-1}{2\ell}}}
    {(2\ell-1)!} \| \nabla w\|_{L^2(\Omega)}^{2\ell-1} \\ 
    & \leq  \| \nabla w \|_{L^2(\Omega)}^3 \sum_{\ell=2}^\infty 
    {\rm TM}(\alpha,\Omega)^{\frac{2\ell-1}{2\ell}}\frac{\alpha^{\frac{2\ell-2}{2}} 2^{4\ell-3} M^{2\ell-1} (\ell!)^{\frac{2\ell-1}{2\ell}}}{(2\ell-1)!}. 
    \end{align*} 
    Once again, one can use the ratio test to show that this last sum converges, which in 
    turn gives us 
    \begin{equation} \label{higher_order6} 
    S_1(w) \lesssim_{\alpha,\Omega,M} \| \nabla w \|_{L^2(\Omega)}^3. 
    \end{equation} 
    Combining \eqref{higher_order4}, \eqref{higher_order5} and \eqref{higher_order6}
    we now see 
    \begin{equation} \label{higher_order7} 
    |T_3(w)| \lesssim_{\alpha, \Omega, M} \| \nabla w \|_{L^2(\Omega)}^3.
    \end{equation} 
    This concludes the proof of the claim.
    
    Finally, by combining \eqref{first_order2}, \eqref{higher_order7}, \eqref{higher_order8}, 
    and \eqref{higher_order9} from Claims 1, 2, 3, respectively, we obtain 
    the desired asymptotic expansion \eqref{eq:taylor_expansion}. 
    The lemma is proved.
\end{proof}

\end{document}
